\section*{Chapitre \ref{ch:b2:climate-variability} --- Variabilité climatique}

\begin{solution}{exo:b2:climate-variability:1}
(a) $29.27 - 26.71 = \qty{2.56}{\celsius}$ en 2026~; $28.79 - 26.73 = \qty{2.06}{\celsius}$ en 1997.
(b) Septembre 2026 était à la fois le mois le plus chaud, de \qty{0.48}{\celsius}, et celui où l'anomalie était la plus forte,
de \qty{0.50}{\celsius}. (c) Chaque année est comparée à une période de référence de 30 ans, centrée sur elle et déplacée
tous les cinq ans~; la \omterm{def:b2:climate-variability:modes}{climatologie} de septembre des années 1980 et 1990 diffère légèrement
de celle des années 2010 et 2020.
\end{solution}

\begin{solution}{exo:b2:climate-variability:2}
(a) JJA~: $(0.79 + 0.95 + 1.28)/3 = 1.007$~; JAS~: $(0.95 + 1.28 + 1.55)/3 = 1.260$~; ASO~: $(1.28 +
1.55 + 1.68)/3 = 1.503$~; SON~: $(1.55 + 1.68 + 2.01)/3 = 1.747$~; OND~: $(1.68 + 2.01 + 2.02)/3 =
1.903$, tous en \unit{\celsius}. (b) Arrondies à une décimale~: 1.0, 1.3, 1.5, 1.7, 1.9, les valeurs de
la table. (c) En ASO 2023, à 1.5.
\end{solution}

\begin{solution}{exo:b2:climate-variability:3}
(a) L'indice suit l'écart de pression entre Tahiti et Darwin, soit la pression de Tahiti moins celle de Darwin~:
inférieure à la moyenne à Tahiti et supérieure à la moyenne à Darwin, il est négatif. (b) El~Niño. (c) Des alizés plus
faibles, voire des vents d'ouest près de la surface à l'ouest, et moins de pluie qu'à l'ordinaire sur l'Indonésie,
car la convection profonde s'est déplacée vers l'est.
\end{solution}

\begin{solution}{exo:b2:climate-variability:4}
(a) $804.1/880.6 = \qty{91.3}{\percent}$ de la moyenne~: dans l'intervalle de 90 à \qty{110}{\percent}, une saison normale,
mais dans sa partie basse. (b) Le total de la saison vaut $556 + 659 + 427 + 252 =
\qty{1894}{mm}$~; juillet en apporte $659/1894 = \qty{34.8}{\percent}$, soit plus d'un tiers.
\end{solution}

\begin{solution}{exo:b2:climate-variability:5}
(a) $\tau = 1025 \times 3992 \times 50/16.1 = \qty{1.271e7}{s} = \qty{147.1}{jours}$ pour
\qty{50}{m}, et le double, \qty{294.2}{jours}, pour \qty{100}{m}. (b) $\mathrm{e}^{-30.44/147.1} =
0.813$ et $\mathrm{e}^{-30.44/294.2} = 0.902$. (c) Le temps d'amortissement est proportionnel à $h$~: une couche de
mélange moins profonde stocke moins de chaleur pour un même flux et oublie une anomalie plus vite~; la mémoire de l'océan,
et la prévisibilité qu'elle procure, diminuent.
\end{solution}

\begin{solution}{exo:b2:climate-variability:6}
(a) Trois mois font 91.3 jours~: $1.2 \times \mathrm{e}^{-91.3/147.1} = \qty{0.65}{K}$~; six mois font
182.6 jours~: $1.2 \times \mathrm{e}^{-182.6/147.1} = \qty{0.35}{K}$. (b) L'anomalie décroît
lentement, c'est donc une meilleure estimation de l'anomalie du mois suivant que zéro. Mais le forçage aléatoire ajoute
chaque jour de nouvelles anomalies, inconnues au moment de la prévision~; à mesure que l'anomalie initiale décroît,
elles prennent une part toujours plus grande, et la prévision dérive vers la \omterm{def:b2:climate-variability:modes}{climatologie}, où elle n'apporte plus rien
d'utile.
\end{solution}

\begin{solution}{exo:b2:climate-variability:7}
(a) $160 \times 86\,400 = \qty{1.38e7}{J}$, soit \qty{13.8}{MJ} par mètre carré et par jour. (b) La
couche contient $(1000 - 850) \times 100/9.81 = \qty{1529}{kg}$ d'air par mètre carré, avec une capacité
thermique de $1529 \times 1004.5 = \qty{1.536e6}{J.m^{-2}.K^{-1}}$~; elle se réchaufferait de $1.38 \times
10^7/1.536 \times 10^6 = \qty{9.0}{K}$ par jour. (c) La couche rayonne sa chaleur dans l'infrarouge, vers l'espace et vers
l'air situé au-dessus~; l'air chaud monte et se mélange à la couche supérieure, ou il est emporté par les vents,
parmi lesquels l'arrivée d'air plus froid que la basse thermique aspire elle-même.
\end{solution}

\begin{solution}{exo:b2:climate-variability:8}
(a) $(-1.93 - 1.11 - 1.98)/3 = -1.673$, soit environ $-1.67$. (b) C'est l'hiver le plus bas de la série.
(c) Un indice NAO fortement négatif~: une dépression d'Islande et un anticyclone des Açores faibles, un courant-jet plus faible
et une trajectoire des tempêtes plus au sud que d'ordinaire~; un hiver froid et sec dans le nord de l'Europe, pluvieux dans le sud de l'Europe,
et davantage d'irruptions d'air froid dans l'est des États-Unis.
\end{solution}

\begin{solution}{exo:b2:climate-variability:9}
(a) $9.3 - 0.4 = \qty{8.9}{mm}$ par jour, plus de 2.5~: la station est dans le domaine. (b) Dans l'hémisphère Sud, l'été
correspond à décembre--février~: $6.1 - 4.2 = \qty{1.9}{mm}$ par jour, moins de 2.5~: la station est hors du domaine,
bien qu'il y pleuve davantage que dans la plupart des points du domaine. (c) $1894/122 = \qty{15.5}{mm}$ par
jour~; Kerala ne sortirait du domaine que si son taux hivernal dépassait $15.5 - 2.5 = \qty{13.0}{mm}$ par jour.
(d) La \omterm{def:b2:climate-variability:monsoon}{mousson} est un renversement saisonnier~: un lieu humide toute l'année, comme la seconde station, n'a pas de
\omterm{def:b2:climate-variability:monsoon}{mousson}, tandis qu'un lieu sec avec une saison des pluies courte et intense en a une.
\end{solution}

\begin{solution}{exo:b2:climate-variability:10}
(a) En dérivant, $\ddot T = a\dot H = -abT$~: l'équation d'un oscillateur harmonique de pulsation
$\omega = \sqrt{ab}$, de période $2\pi/\sqrt{ab}$. (b) $2\pi/\sqrt{ab} = 4$ donne $ab =
(2\pi/4)^2 = \qty{2.47}{an^{-2}}$. (c) Avec $T = u\,\mathrm{e}^{-rt}$, $H = v\,\mathrm{e}^{-rt}$,
$\dot T = (\dot u - ru)\,\mathrm{e}^{-rt} = (av - ru)\,\mathrm{e}^{-rt}$, donc $\dot u = av$, et de même
$\dot v = -bu$~: $(u, v)$ oscille sans amortissement, et $(T, H)$ est cette oscillation multipliée par
$\mathrm{e}^{-rt}$. En une période, l'amplitude est multipliée par $\mathrm{e}^{-0.25 \times
4} = \mathrm{e}^{-1} = 0.37$. (d) Un forçage~: les fluctuations aléatoires des alizés, qui frappent le système
encore et encore, et la \omterm{def:b2:climate-variability:bjerknes}{rétroaction de Bjerknes}, qui rend l'amortissement plus faible,
voire négatif.
\end{solution}

\begin{solution}{exo:b2:climate-variability:11}
(a) $A = R^2 \,\Delta\lambda\,(\sin\ang{5} - \sin(\ang{-5})) = (6.371 \times 10^6)^2 \times 0.8727
\times 0.1743 = \qty{6.17e12}{m^2}$. (b) $Q = 1025 \times 3992 \times 100 \times 2 \times 6.17
\times 10^{12} = \qty{5.05e21}{J}$. (c) Six mois représentent $182.6 \times 86\,400 = \qty{1.578e7}{s}$~;
par mètre carré, $8.18 \times 10^8/1.578 \times 10^7 = \qty{51.9}{W/m^2}$, maintenu pendant un semestre.
(d) Depuis l'océan lui-même~: des alizés plus faibles remontent moins d'eau froide à l'est, et l'eau chaude de l'ouest,
stockée au-dessus d'une thermocline plus profonde, s'étend vers l'est~; une grande part du réchauffement
peut provenir de la chaleur déjà présente dans l'océan, déplacée plutôt que gagnée.
\end{solution}

\begin{solution}{exo:b2:climate-variability:12}
(a) La boucle restitue $0.3g$, puis $0.3g^2$, et ainsi de suite~: le total vaut $0.3(1 + g + g^2 + \dots) =
0.3/(1 - g)$, série géométrique qui converge pour $g < 1$~: \qty{0.6}{K}, \qty{1.5}{K} et
\qty{3.0}{K} pour $g = 0.5$, 0.8 et 0.9. (b) La série diverge~: l'anomalie croîtrait sans limite~; la croissance doit
être arrêtée par des processus que la boucle omet, l'amortissement par le flux de chaleur de surface et la vidange
du stock d'eau chaude que décrit l'oscillateur de recharge. (c) Plus la surface de la mer est chaude, plus un
léger réchauffement déclenche facilement une convection profonde et des pluies, dont le chauffage entretient les vents~;
c'est la saison où les effets de l'ENSO sont les plus forts.
\end{solution}

\begin{solution}{pb:b2:climate-variability:1}
\textbf{1.} $(0.26 + 0.75 + 1.13)/3 = \qty{0.713}{\celsius}$.
\textbf{2.} MJJ $(0.75 + 1.13 + 1.52)/3 = 1.133$~; JJA $(1.13 + 1.52 + 1.80)/3 = 1.483$~; JAS $(1.52 +
1.80 + 2.06)/3 = 1.793$.
\textbf{3.} ASO $(1.80 + 2.06 + 2.26)/3 = 2.040$~; SON $(2.06 + 2.26 + 2.42)/3 = 2.247$~; OND $(2.26
+ 2.42 + 2.36)/3 = 2.347$~; NDJ $(2.42 + 2.36 + 2.33)/3 = 2.370$.
\textbf{4.} Arrondies à une décimale, 0.7, 1.1, 1.5, 1.8, 2.0, 2.2, 2.3 et 2.4~: la table.
\textbf{5.} Un seul mois porte du bruit météorologique~; l'ENSO dure plusieurs mois, et la moyenne de trois mois
élimine une grande part du bruit tout en conservant le signal lent.
\textbf{6.} En JJA 1997, à 1.483, arrondi à 1.5~; il culmine à 2.4 en NDJ.
\textbf{7.} De ASO 2018 à MJJ 2019, et de SON 2019 à FMA 2020.
\textbf{8.} Dix saisons, et six.
\textbf{9.} En ASO 2020~; au sein de 2020, seules quatre saisons, de ASO à NDJ, sont à ou au-dessous de $-0.5$, il en manque
une pour la règle, et DJF 2021 la complète.
\textbf{10.} Un El~Niño, de SON 2018 à MAM 2019, sept saisons (AMJ 2019 tombe à 0.37)~; une
La~Niña commence en AMJ 2020 ($-0.55$), quatre saisons avant la place que lui donne l'ONI, en ASO.
\textbf{11.} Non~: l'indice relatif ne dépasse jamais 0.29 à la fin de 2019. Tout l'océan tropical était alors
plus chaud que sa période de référence d'environ un demi-degré~; la région Niño-3.4 était chaude en valeur absolue,
mais à peine plus chaude que les tropiques qui l'entourent.
\textbf{12.} Onze.
\textbf{13.} 8, 7, 10, 4, 5, 6, 12, 5, 9 et 3 ans.
\textbf{14.} Le plus court est de 3 ans (2023 à 2026), le plus long de 12 (1997 à 2009).
\textbf{15.} $69/10 = 6.9$ ans dans les deux cas~: la somme des intervalles successifs est la durée entre le
premier début et le dernier, $2026 - 1957 = 69$ ans.
\textbf{16.} Il dure déjà depuis cinq saisons, de MAM à JAS, à \qty{0.5}{\celsius} au moins, et il a culminé à 2.2,
au-dessus de 1.5~; son début, MAM 2026, est la première saison de la série, et les saisons suivantes ne peuvent
qu'allonger la série, sans en déplacer le commencement.
\textbf{17.} Non, et ce n'est pas nécessaire~: les 25 épisodes de toutes forces arrivent environ tous les trois ans, comme l'indique la FAQ~;
les forts n'en sont qu'un sous-ensemble, donc plus rares, et leurs intervalles vont de 3 à 12 ans.
\textbf{18.} $\tau = 1025 \times 3992 \times 50/16.1 = \qty{1.27e7}{s} = \qty{147}{jours}$.
\textbf{19.} Couche de mélange~: $1025 \times 3992 \times 50 = \qty{2.05e8}{J.m^{-2}.K^{-1}}$~; colonne d'air~:
$(10^5/9.81) \times 1004.5 = 10\,194 \times 1004.5 = \qty{1.02e7}{J.m^{-2}.K^{-1}}$, 20 fois
moins.
\textbf{20.} $1.024 \times 10^7/(1025 \times 3992) = \qty{2.5}{m}$ d'eau de mer.
\textbf{21.} $1.024 \times 10^7/16.1 = \qty{6.36e5}{s} = 7.4$~jours.
\textbf{22.} Une anomalie de l'air seul s'amortit en une semaine environ, celle de la couche de mélange en
cinq mois environ, et la couche chaude au-dessus de la thermocline est plus profonde encore~: seul l'océan retient
la chaleur assez longtemps pour porter une anomalie d'une saison à la suivante, tandis que l'atmosphère fournit les
vents qui couplent les deux.
\textbf{23.} Depuis 1950, les épisodes forts d'El~Niño ont commencé en moyenne tous les \textbf{6.9 ans},
soit environ sept.
\end{solution}
